\documentclass[11pt]{article}
\usepackage[margin=1in]{geometry}
\usepackage{amsmath,amssymb,amsthm}
\usepackage[hidelinks]{hyperref}
\newtheorem{theorem}{Theorem}
\newtheorem{lemma}{Lemma}
\title{A proof of the generalized power-sum identity\\
\large MathDB 334350, with an extension to every odd function}
\author{Principia Math In Progress}
\date{18 September 2026}
\begin{document}
\maketitle
\begin{abstract}
We prove Svinin's conjectured equality between a sum over weakly
increasing tuples and a binomial combination of sums over ordered
tuples. In fact, the equality holds when the odd power is replaced
by any odd function on the integers. The proof establishes a
Laurent polynomial identity by counting gaps and applying an
elementary truncated negative-binomial identity.
\end{abstract}

\section*{Task 1: the original conjecture and its full solution}
Let $n,j\geq1$ be integers, and define
\[
 A(z)=z+\cdots+z^n,\qquad
 A(z)^\ell=\sum_{q=0}^{\ell(n-1)}C_{\ell,q}(n)z^{\ell+q}.
\]
For $k\geq0$, put
\[
 s_{k,\ell}(n)=\sum_{q=0}^{\ell(n-1)}C_{\ell,q}(n)(\ell+q)^{2k+1},
 \qquad
 S_{k,j}(n)=\sum_{q=0}^{j-1}\binom{j(n+1)}q s_{k,j-q}(n).
\]
Conjecture 3.1 of \cite{S} asserts that
\begin{equation}\label{target}
 \sum_{1\leq\lambda_1\leq\cdots\leq\lambda_j\leq jn}
       \sum_{i=1}^j(\lambda_i-(i-1)n)^{2k+1}=S_{k,j}(n).
\end{equation}
The source proves the case $j=2$. We prove every $j$, without
restrictions on $k$ or $n$ beyond those stated.

\begin{theorem}\label{main}
For every odd function $\psi:\mathbb Z\to\mathbb C$ and all
$n,j\geq1$,
\begin{equation}\label{strong}
 \sum_{1\leq\lambda_1\leq\cdots\leq\lambda_j\leq jn}
       \sum_{i=1}^j\psi(\lambda_i-(i-1)n)
 =\sum_{\ell=1}^j\binom{j(n+1)}{j-\ell}
       \sum_{r_1,\ldots,r_\ell=1}^n\psi(r_1+\cdots+r_\ell).
\end{equation}
In particular, \eqref{target} holds for every $k\geq0$.
\end{theorem}

\subsection*{A finite coefficient identity}
\begin{lemma}\label{truncation}
For positive integers $a,j$, define
\[
 P(X)=\sum_{q=0}^{j-1}\binom{a+q-1}q X^{j-q},\qquad
 T(b)=[t^{a-1}]\frac1{(1-t)^j(b-t)}.
\]
As rational functions of $b$,
\begin{equation}\label{T}
 T(b)=b^{-a}(1-b)^{-j}-P\bigl((1-b)^{-1}\bigr).
\end{equation}
\end{lemma}
\begin{proof}
The elementary polynomial identity
\begin{equation}\label{binomial}
 (1-b)^j\sum_{p=0}^{a-1}\binom{j+p-1}p b^p
 +b^a\sum_{q=0}^{j-1}\binom{a+q-1}q(1-b)^q=1
\end{equation}
follows by partitioning binary words of length $a+j-1$ according
to whether the $j$th second symbol or the $a$th first symbol
occurs first. Give the first symbol weight $b$ and the second
weight $1-b$. In the first case, $p<a$ first symbols precede the
last second symbol; in the second case, $q<j$ second symbols
precede the last first symbol. The total weight of all possible
suffixes is one. These two cases give the two sums, and the total
weight of all words is $(b+(1-b))^{a+j-1}=1$.

On the other hand, expanding at $t=0$ gives
$T(b)=b^{-a}\sum_{p=0}^{a-1}\binom{j+p-1}p b^p$.
Divide \eqref{binomial} by $b^a(1-b)^j$ to obtain \eqref{T}.
\end{proof}

\subsection*{The marked-tuple generating polynomial}
Fix $n,j$, and abbreviate
\[
 a=jn,\qquad h=n+1,\qquad N=j(n+1)=a+j,\qquad
 B(z)=1+A(z).
\]
Let
\[
 F(z)=\sum_{1\leq\lambda_1\leq\cdots\leq\lambda_j\leq a}
             \sum_{i=1}^j z^{\lambda_i-(i-1)n},\qquad
 R(z)=\sum_{\ell=1}^j\binom N{j-\ell}A(z)^\ell.
\]
Thus $F$ is a Laurent polynomial and $R$ is an ordinary
polynomial with zero constant term. We will prove
\begin{equation}\label{Laurent}
 F(z)-F(z^{-1})=R(z)-R(z^{-1}).
\end{equation}

The involution
$\lambda_i\mapsto a+1-\lambda_{j+1-i}$ sends each exponent
to $h$ minus the corresponding reversed exponent. Therefore
\begin{equation}\label{symmetry}
 F(z^{-1})=z^{-h}F(z).
\end{equation}
For a tuple $\lambda$, write its nonnegative gaps as
\[
 g_0=\lambda_1-1,\quad
 g_i=\lambda_{i+1}-\lambda_i\ (1\leq i<j),\quad
 g_j=a-\lambda_j.
\]
They sum to $a-1$, and
$\lambda_i-(i-1)n=1+g_0+\cdots+g_{i-1}-(i-1)n$.
Summing each gap independently gives
\begin{align}
 F(z)&=z[t^{a-1}]\sum_{l=0}^{j-1}
 \frac{z^{-nl}}{(1-zt)^{l+1}(1-t)^{j-l}}\label{gaps}\\
 &=z[t^{a-1}]
 \frac{(1-t)^{-j}-z^{-a}(1-zt)^{-j}}
      {(1-zt)-z^{-n}(1-t)}.\label{geometric}
\end{align}
The second equality is a finite geometric sum.

Work temporarily in the field $\mathbb Q(z)$ and put
\[
 b=\frac{z^n-1}{z^{n+1}-1}.
\]
Multiplying \eqref{geometric} by $1-z^{-h}$ and using
\eqref{symmetry} yields
\begin{align}
 F(z)-F(z^{-1})
 &=[t^{a-1}]\frac{(1-t)^{-j}-z^{-a}(1-zt)^{-j}}{b-t}
 \notag\\
 &=T(b)-T(zb).\label{difference}
\end{align}
For the second term substitute $w=zt$ in the coefficient
extraction; the factors $z^{a-1}$, $z$, and $z^{-a}$ cancel.

Now
\[
 1-b=z^n(1-zb),\qquad
 (1-zb)^{-1}=B(z),\qquad
 (1-b)^{-1}=z^{-n}B(z)=B(z^{-1}).
\]
Since $a=jn$, we also have
\[
 b^{-a}(1-b)^{-j}=(zb)^{-a}(1-zb)^{-j}.
\]
These terms cancel when Lemma \ref{truncation} is applied to
\eqref{difference}, leaving
\begin{equation}\label{Pdifference}
 F(z)-F(z^{-1})=P(B(z))-P(B(z^{-1})).
\end{equation}

\subsection*{Identifying the binomial combination}
As an identity of polynomials in $B$, with $A=B-1$, we have
\begin{align*}
 R+\binom Nj
 &=[w^j]\frac{(1+w)^N}{1-(B-1)w}\\
 &=[w^j]\sum_{\ell\geq0}B^\ell w^\ell(1+w)^{N-1-\ell}\\
 &=\sum_{\ell=0}^j\binom{N-1-\ell}{j-\ell}B^\ell
  =P(B)+\binom{N-1}j.
\end{align*}
Thus $R=P(B)-\binom{N-1}{j-1}$. The constant disappears upon
subtracting the expression with $z^{-1}$, so
\eqref{Pdifference} proves \eqref{Laurent}. Although rational
functions were used to establish it, both sides of
\eqref{Laurent} are Laurent polynomials. There are no exceptional
values of $n,j$ or numerical substitutions to exclude.

\begin{proof}[Proof of Theorem \ref{main}]
Apply to \eqref{Laurent} the linear map sending $z^r$ to
$\psi(r)$. Oddness gives $\psi(-r)=-\psi(r)$ and $\psi(0)=0$,
so each side becomes twice its evaluation on $F$ or $R$,
respectively. Dividing by two proves \eqref{strong}. Expanding
$A(z)^\ell$ and choosing $\psi(r)=r^{2k+1}$ gives exactly the
source coefficients and \eqref{target}.
\end{proof}

As a further consequence, if $c_r$ is the number of marked tuples
with exponent $r$ minus the number with exponent $-r$, then
\[
 c_r=[z^r]R(z)>0\qquad(1\leq r\leq jn).
\]
Indeed, choose $\ell=\lceil r/n\rceil\leq j$; since
$\ell\leq r\leq\ell n$, at least one ordered $\ell$-tuple from
$\{1,\ldots,n\}$ has sum $r$. This proves the positivity
consequence noted conditionally in \cite{SS}.

\section*{Task 2: a direct cancellation on marked tuples}
Construct an explicit sign-reversing pairing between marked
tuples with negative and positive exponents. Identify the
unpaired positive objects with the ordered tuples and binomial
choices on the right of \eqref{strong}. Such a construction would
give a direct combinatorial explanation of the coefficient
positivity proved above.

\section*{Task 3: the even part of the distribution}
Determine a useful positive formula for $F(z)+F(z^{-1})$.
Use it to express sums of even functions over the same weak
tuples and to quantify the difference from the odd-function
binomial transform. Include the mass at exponent zero.

\section*{Task 4: quantitative properties of the positive coefficients}
Study the coefficient sequence $([z^r]R(z))_{r=1}^{jn}$.
Determine its modes, whether it is unimodal or log-concave, and
its limiting distribution when $j$ or $n$ grows. The strict
positivity just proved is a starting point, not a solution to
these further questions.

\section*{Task 5: a recurrence for the normalized power sums}
Use the identity to investigate the source's separate Conjecture
5.1 on representing $S_{k,j}(n)$ by polynomials in $j$ with
coefficients rational in $n(n+1)$. Seek a recurrence uniform in
$n$, while distinguishing the known special case $n=1$ and the
later partial results cited in the literature audit.

\paragraph{Verification and scope.}
Exact tuple counts, independent convolutions, and rational checks
supplement the all-parameter proof. Tasks 2--5 are proposals.
Review is by the authoring agent; no independent review or proof
assistant is claimed.

\paragraph{Campaign eligibility.}
This is a full mathematical proof with zero new campaign credit
pending the literature audit. A related 2019 paper concerning the
same sums was accessible only as an abstract and first-page
preview. Its remaining contents must be checked before eligibility
can be cleared; see the accompanying literature audit.

\begin{thebibliography}{9}
\bibitem{S} A. K. Svinin, \emph{Conjectures Involving a Generalization
of the Sums of Powers of Integers}, Experimental Mathematics
\textbf{27}(4), 437--443; published online 5 April 2017,
Conjecture 3.1. \url{https://doi.org/10.1080/10586458.2017.1306815}.
Preprint: arXiv:1610.05387v2 (1 March 2017).
\bibitem{SS} A. K. Svinin and S. V. Svinina,
\emph{On Some Generalizations of Sum of Powers of Natural Numbers},
Russian Mathematics \textbf{63}(8) (2019), 25--37.
\url{https://doi.org/10.3103/S1066369X19080048}.
The Russian original, pages 37--38, equation (25), retains the
target as a conjecture and states its conditional positivity
consequence.
\end{thebibliography}
\end{document}
